\scanpage{50}

% 1–3
Un $\Sigma$-monomorphisme (épimorphisme) $A\longrightarrow B$ de $B$ est un morph. qui s'insère dans une suite $\in\Sigma$

% 4
\[
 0\longrightarrow A\longrightarrow B\longrightarrow K\longrightarrow0\quad(\text{resp.}\quad0\longrightarrow P\longrightarrow A\longrightarrow B\longrightarrow0)
\]

% 5–7
Un morphisme \inserted{$A\longrightarrow B$} \deleted{\unsure{est}} est \inserted{$\Sigma$-}strict, s'il a un noyau $P$ et un conoyau $Q$, \deleted{\unsure{et si|objet} $[\operatorname{Coim}\overset{\sim}{\longrightarrow}\operatorname{Im}$,} \inserted{et si $P\longrightarrow A$} $(B\longrightarrow Q)$ est un $\Sigma$-mono ($\Sigma$-épi),

% 5–7, outer left margin
% slanted
\margin{définir \unsure{$\Sigma$-mo.|$\Sigma$-simpl.} et \unsure{$\Sigma$-épi|$\Sigma$-co.} \deleted{\illegible{1}}}

% 7–9
Il revient au même de dire que l'on peut trouver des suites $\in\Sigma$

% 10–11
\[
 \begin{aligned}
 0&\longrightarrow P\longrightarrow A\longrightarrow R\longrightarrow0\\
 0&\longrightarrow R\longrightarrow B\longrightarrow Q\longrightarrow0
 \end{aligned}
\]

% 12–14
telles que $A\longrightarrow B$ soit le composé $A\to R\longrightarrow B$. (Il faut peut-être exiger qu'un composé de $\Sigma$-mono (resp $\Sigma$-épi) est itou.)

% 15–16
On définit de façon évidente les \inserted{objets} $\Sigma$-projectifs et $\Sigma$-injectifs. Une \deleted{\inserted{suite} complexe est dit} \deleted{complexe de cochaînes (disons)}

% 17–18, cancelled
\deleted{$\longrightarrow C^{i-1}\longrightarrow \overset{\deleted{\unsure{d}}}{C^i}\longrightarrow C^{i+1}\longrightarrow\cdots$ est dit spécial.}

% 19–22
\[
 A\xrightarrow{u}B\xrightarrow{v}C
\]
\deleted{est dite spéciale si \unsure{$vu=0$}, et} $\Sigma$-spéciale \deleted{\unsure{si}} si elle \deleted{elle \unsure{déf.} s'insère dans des suites (de $\Sigma$)} \deleted{$vu=0$, $u$ est spécial et s'insère ds une suite \unsure{$v=0$|$vu=0$}, \unsure{sa}} de suites exactes $\in\Sigma$

% 23–26, stepped sequences
\diagram{\begin{tikzcd}[ampersand replacement=\&,column sep=.45em,row sep=1.3em]
 0\arrow[r]\&\operatorname{Ker}u\arrow[r]\& A\arrow[r]\&\operatorname{Coim}u\arrow[r]\arrow[d,"\wr"]\&0\&\&\&\&\&\\
 \&\&0\arrow[r]\&\operatorname{Im}u\arrow[r]\&\operatorname{Ker}v\arrow[r]\arrow[d,equal,"\wr"]\&H\arrow[r]\&0\&\&\&\\
 \&\&\&0\arrow[r]\&\operatorname{Ker}v\arrow[r]\& B\arrow[r]\&\operatorname{Coim}v\arrow[r]\arrow[d,"\wr"]\&0\&\&\\
 \&\&\&\&\&\deleted{0}\arrow[r]\&\operatorname{Im}v\arrow[r]\& C\arrow[r]\&\unsure{\operatorname{Cok}v|\operatorname{Coker}v}\arrow[d]\&\\
 \&\&\&\&\&\&\&\&0\&
\end{tikzcd}}

% 27–32
(c'est une notion autoduale) On dit qu'un complexe $C^{\bullet}$ dans $B$ est \uline{$\Sigma$-spécial} si …, $\Sigma$-acyclique si de \uline{plus} les $H^i(C^{\bullet})$ sont nuls, une $\Sigma$-résolution si c'est un complexe de cochaînes ($C^i=0$ si $i<0$) et $H^i=0$ pour $i\deleted{\unsure{\geq}}\unsure{>|\neq|\geq}0$.

% archive, bottom left
% archive sheet number: 49
